\section{How Template Competition Opens a Scale Band}\label{sec:geometry}
Fix a state $(s,t)$ with $0<s<1/2$ in the three-template geometry
and take $v_L>\tau^2$.
For inverse variance $z=1/v$ and $z_L=1/v_L$, define the half squared-distance
gaps relative to template 0:
\begin{equation}\label{eq:band_gaps}
 P=\tfrac12-s,\qquad Q=\tfrac12(a^2+d^2+2as-2dt).
\end{equation}
The responsibility odds relative to template 0 are $r e^{-Pz}$ and
$h e^{-Qz}$. Write $m(z)$ for the first coordinate of the posterior
template mean $\bar\mu_v$. Increasing inverse variance from zero suppresses template 2 faster than template 1 if $Q>P$.
We give the precise threshold for that suppression to open a band.

\begin{proposition}[Exact mass window at a fixed state]\label{prop:mass_band}
Fix a posterior-mean target $\rho\in(0,1)$ and define
\begin{equation}\label{eq:mass_capacity}
 H_\rho(z)=
 \frac{(1-\rho)r\,e^{(Q-P)z}-\rho e^{Qz}}{a+\rho}.
\end{equation}
If $Q\le P$, an interior variance cannot attain $m(z)>\rho$ when
both endpoints fail to attain it. If $Q>P$, set
\begin{equation}\label{eq:band_peak}
 z_\rho=\left[\frac{1}{P}
     \log\frac{(Q-P)(1-\rho)r}{Q\rho}\right]_{[0,z_L]},
\end{equation}
where the brackets clip to the specified interval.
Both endpoints have $m\le\rho$ and some intermediate variance has
$m>\rho$ exactly when
\begin{equation}\label{eq:exact_mass_window}
 \max\{H_\rho(0),H_\rho(z_L)\}\le h < H_\rho(z_\rho).
\end{equation}
The working band is the interval determined by $H_\rho(z)>h$, mapped
back to variance by $v=1/z$.
\end{proposition}

Indeed, $m(z)>\rho$ is equivalent to
\[
 (1-\rho)r>\rho e^{Pz}+(a+\rho)h e^{-(Q-P)z}.
\]
The right side is nondecreasing if $Q\le P$ and strictly convex if
$Q>P$. For the latter, the sign of $H_\rho'$ is the sign of
$(Q-P)(1-\rho)r-Q\rho e^{Pz}$. Thus the upper mass threshold
and the peak in~\eqref{eq:band_peak} are explicit. The lower threshold
keeps the two endpoints below the target. For fixed $a$ and other
inputs, raising $h$ shrinks the working interval, while increasing
$Q$ expands it by weakening the competition at each $z>0$.
At $\rho=s$, the result exactly detects whether the prior part
of the RED update can point right when both endpoints fail to do so.
It applies to positive Gaussian component widths.

\subsection{From the mass window to an update threshold}
A desired horizontal displacement per unit step size $d_*$ gives
the normalized prior target $u=(d_*-b_1+\alpha s)/\lambda$.
Set $l=1-\tau^2 z$, $A_u=(1-s)l-u$ and
\begin{equation}\label{eq:update_ratios}
 U(z)=\frac{s l+u}{A_u},\qquad
 W(z)=\frac{(a+s)l+u}{A_u}.
\end{equation}
The exact RED displacement per unit step exceeds $d_*$ precisely when
\begin{equation}\label{eq:update_band}
 J_u(z):=U(z)e^{Pz}+hW(z)e^{-(Q-P)z}<r
\end{equation}
where $A_u>0$. This uses the full $c(v)$ and the actual data force.
For $u=0$, it recovers Proposition~\ref{prop:mass_band} at $\rho=s$.

\begin{proposition}[A convex test for the RED variance band]\label{prop:convex_band}
Let $Q>P>0$, $v_L>\tau^2$, and assume $A_u(z_L)>0$ and
$s l(z_L)+u>0$. If $u\ge0$, $J_u$ is strictly convex on $[0,z_L]$.
For $u<0$ the same conclusion holds if
\begin{align}
 P^2U_L&>2P d_1+d_2,\nonumber\\
 (Q-P)^2 W_L&>(a+1)d_2,\label{eq:band_convexity}
\end{align}
where $U_L,W_L,A_L$ denote the terminal values of $U,W,A_u$, and
\[
 d_1=|u|\tau^2/A_L^2,\qquad
 d_2=2(1-s)\tau^2 d_1/A_L.
\]
Under these conditions, a nonempty band with both endpoints failing
exists exactly when $\min J_u<r\le\min\{J_u(0),J_u(z_L)\}$.
Its endpoints are the level roots of $J_u=r$.
\end{proposition}

The minimum is obtained by clipping the zero of the increasing
derivative $J_u'$ to the inverse-variance interval, or by taking the
appropriate endpoint if no zero exists. For a test of \emph{endpoint
dominance itself}, choose $d_*$ as the larger endpoint displacement
per unit step size. Whenever the hypotheses of
Proposition~\ref{prop:convex_band} hold at that target, its test
decides whether any intermediate step exceeds both endpoints.
The convexity conditions are checked at the chosen target.

\paragraph{Additional local thresholds.}
Supplement S7 derives an explicit critical competitor mass for
improvement over the endpoint steps, and S24 treats several
competitors. The risk family below uses the full drift test
\eqref{eq:update_band}, which also applies outside that critical-mass
subfamily.

\paragraph{Safe scales for a candidate barrier.}
Consider $s\le\xi,\ t\ge-\eta_B$ with $0<\xi<1/2$.
At $(\xi,-\eta_B)$ use $d_*=-\delta/\gamma$, for coordinate
update errors bounded by $\delta$. Provided the lower face is
inward and each map is nondecreasing in its own coordinate,
the variances with $J_u\ge r$ preserve the region. The other
variances admit an outward update at its corner under the allowed
error bound. Thus the roots predict where the common barrier loses
invariance for a given observation. To certify a schedule, a
separate channel argument checks progress and transverse
invariance along the trajectory, as in Theorem~\ref{thm:ind_noise}.

At the same state, the nominal recovery-box mass $h=1.78$ and
noiseless class-1 data with $\alpha=10^{-3}$ give, at
$\lambda=1+\alpha^2$ and $\delta=10^{-4}$, a corner-outward
variance interval of approximately $(.18064,1.87551)$.
The interval comes from the full update test~\eqref{eq:update_band}.
With $\delta=0$, the positive-drift interval for the exact update
is approximately $(.18095,1.87272)$. The larger interval above
expresses failure of robust invariance. At each of its scales some
allowed perturbation points outward from the corner. Reaching that
corner and entering a correct terminal rectangle require separate
trajectory arguments.
The restricted union lies outside both intervals, and the certified
working band $[.69,.71]$ lies inside. This calculation concerns a
single model and the candidate barrier. The uniform risk theorem
uses separate bounds across the parameter box and the channel.
